\begin{figure}[tb]
\centering
\begin{tikzpicture}[x=1cm,y=1cm,>=Stealth,font=\small]
  \fill[blue!12] (0,0) rectangle (1.2,.55);
  \fill[gray!7] (1.2,0) rectangle (8.8,.55);
  \fill[orange!18] (8.8,0) rectangle (10,.55);
  \draw[thick] (0,0) -- (10,0);
  \foreach \x in {0,1.2,8.8,10} {\draw (\x,-.08) -- (\x,.62);}
  \node[below] at (0,-.09) {$b_-$};
  \node[below] at (1.2,-.09) {$b_-+\delta$};
  \node[below] at (8.8,-.09) {$b_+-\delta$};
  \node[below] at (10,-.09) {$b_+$};
  \node at (.6,.3) {$r_n=0$};
  \node at (9.4,.3) {$r_n=1$};
  \node at (5,.3) {either choice may occur};
  \draw[->,thick,blue!65!black] (.12,.9) -- (1.1,.9);
  \node[above,blue!65!black] at (.61,.97) {$+\delta t$};
  \draw[->,thick,orange!70!black] (9.88,.9) -- (8.9,.9);
  \node[above,orange!70!black] at (9.39,.97) {$-\delta(1-t)$};
\end{tikzpicture}
\caption{The bounded recurrence, schematically. The tested bit forces
an inward step in each endpoint strip; the bit need not be monotone
between them. Each step has length less than $\delta$, so no state
can leave $[b_-,b_+]$. The oracle computation produces the choices
$r_n$, not the numerical states $s_n$.}
\label{fig:bounded-recurrence}
\end{figure}
